\documentclass[12pt]{report}
\usepackage{cite}
\usepackage{verbatim}
\usepackage[font=small,labelfont=bf]{caption}
\usepackage{graphicx}
\usepackage{fullpage}
%\usepackage{soul}
\usepackage[german]{babel}
\usepackage{float}
\usepackage{mathtools}
\usepackage{amsfonts}
%\usepackage{hyperref}
%\usepackage{amssymb}
%\usepackage{paralist}
%\usepackage{fancyhdr}
\usepackage{amsthm}
\usepackage{setspace}
\usepackage{subfig}
\usepackage{listings}
\usepackage{xspace}
\usepackage{listings}
\usepackage{fontawesome}
\usepackage{amsmath}
\usepackage{amssymb}
\usepackage{xcolor}
\usepackage{soul}
\newtheorem{theorem}{Satz}
\newtheorem{lemma}[theorem]{Lemma}
\newtheorem{Beweis}[theorem]{Beweis}
\newtheorem{Korollar}[theorem]{Korollar}
\newtheorem{Beispiel}[theorem]{Beispiel}
\newtheorem*{Bemerkung*}{Bemerkung}
\newtheorem*{Konstruktion*}{Konstruktion}
\newtheorem{Definition}[theorem]{Definition}
\newcommand{\PsubT}{\text{Principal subtree }}
\newcommand{\linpreft}{\text{linear preferential attachment trees }}



% Title Page
% LTeX: language=de-DE
\begin{document}
\author{Victor Velesco}

\begin{titlepage}
\vspace*{1cm}
\begin{center}
\large
\textbf{Johann Wolfgang Goethe-Universität Frankfurt am Main}\\
Fachbereich 12 Mathematik/ Institut für Mathematik\\
Studiengang Mathematik (Bsc.)\\
\vspace{1cm}
\Large
\textbf{Bachelorarbeit im Bereich Stochastische Analyse von Algorithmen}\\
\vspace{0.5cm}
\Large
\textit{Preferential attachment Trees}\\
\vspace{3cm}
\includegraphics[width = 4cm]{Pictures/1280px-Goethe-Logo.svg.png}
\vspace{3cm}

\end{center}

Victor Velesco\\
Matrikelnummer: 7419816\\
Fachsemester: 8\\ 
Betreuender Professor: Prof. Dr. Ralph Neininger\\
Zweitgutachter: Jasper Ischebeck\\
\end{titlepage}
\tableofcontents
\chapter{Einführung}
% LTeX: language=de-DE
Die vorliegende Bachelorarbeit ist eine Ausarbeitung von der wissenschaftlichen Publikation \glqq Random Recursive Trees and Preferential Attachment Trees are Random Split Trees\grqq, von Svante Janson \cite{janson2019random}.\\
Die Publikation von Janson verallgemeinert die von Devroye \cite{devroye1998universal} eingeführten \textit{Split trees}, indem auch Bäume mit unendlichen Ausgangsgraden, also Anzahl von Kindern pro Wurzel, zugelassen werden. Präziser formuliert, werden sogenannte \textit{Linear Preferential Attachment Trees} in die Vorlage der Split-Bäume aufgenommen. Je nach Parameter Wahl umfassen die \textit{Preferential Trees} unter anderem die \textit{m-ary, standard preferential und zufälligen rekursiven} Bäume und Eigenschaften der jeweiligen Bäume werden in \textcolor{red}{Füge später hinzu} näher betrachtet. Da die Publikation stark auf \textit{Kingman's theory of exchangeable partitions} \cite{kingman1978representation,kingman1982coalescent} als auch auf \textit{Pitman's chinese restaurant process} \cite{pitman2006combinatorial,pitman1995exchangeable} beruht, wird im ersten Kapitel ein Überblick über diese Themengebiete verschaffen. Anschließend wird das Haupttheorem zur Einbettung in Kapitel 2 bewiesen, gefolgt von einer Anwendung zu Tiefen in Kapitel 3. Im letzten Kapitel wird ein Spezialfall der \textit{Preferential Trees} gesondert behandelt.

\begin{section}{Motivation}
% LTeX: language=de-DE
Als vorläufige Motivation für die Relevanz der Theorie der austauschbaren und partiell austauschbaren Partitionen von $\mathbb{N}$ für die darauffolgende Baumanalyse, beschreiben wir kurz informell den Zusammenhang. Die Konstruktion von Bäumen die betrachtet werden, beinhalten sukzessives Hinzufügen von Knoten, welche mit natürlichen Zahlen identifiziert werden. Teilbäume, wo die Knoten/natürliche Zahlen im Baum angefügt werden, kann man sich deshalb als eine Art Partition vorstellen. Die Austauschbarkeit erlaubt uns eine De-Finetti-ähnlichen Aussage (Satz \ref{Main theorem CRP}) über die relative Größe von Partitionen und somit über die Größe von Teilbäumen zu formulieren. \\
Wir beginnen mit den Definitionen. 
\end{section}
\begin{section}{Definitionen}
\begin{Definition}
   Sei $F$ eine endliche Menge. Dann ist eine Partition von $F$ eine ungeordnete Sammlung von nichtleeren, disjunkten Untermengen $\{A_i| 1 \leq i \leq k, A_i \subseteq F\}$, mit $\cup_{i=1}^{k}A_k = F$.\\
   Eine Partition von einer natürlichen Zahl $n$, ist eine ungeordnete Menge von positiven natürlichen Zahlen mit Summe $n$.
\end{Definition}
\begin{Definition}
    \textnormal{(Zufällige Partitionen)}
\begin{enumerate}
    \item Eine zufällige Partition von $n\in \mathbb{N}$ ist eine Zufallsvariable $\sigma_n$ mit Werten in der endlichen Menge von allen Partitionen von $n$.
    \item Eine zufällige Partition von $[n] := \{1 ...n\}$, ist eine Zufallsvariable $\Pi_n$ mit Werten in $\mathcal{P}(n)$, der endlichen Menge von Partitionen von $[n]$.
    \item Eine zufällige Partition auf $\mathbb{N}$ ist eine Folge $\Pi = (\Pi_n)_{n\in \mathbb{N}}$ von Zufallsvariablen auf einem gemeinsamen Wahrscheinlichkeitsraum mit Werten in der endlichen Menge von Partitionen von $[n]$, sodass für $i<j, i,j \in \mathbb{N}$ die Restriktion von $\Pi_j$ auf $[i]$ $\Pi_i$ ist.
\end{enumerate}

\end{Definition}
\begin{Definition}
    \textnormal{(Austauschbarkeit)}
    \label{Austauschbarkeit}
    \begin{enumerate}
        \item Eine zufällige Partition $\Pi_n$ von $[n]$ heißt partiell Austauschbar, falls für jede Partition $\{ A_1 ... A_k\}$ von $[n]$ in Reihenfolge ihres Erscheinens\footnote{\textit{In Reihenfolge ihres Erscheinens} bedeutet $1 \in A_1$, und für alle $2 \leq i \leq k$, das erste Element von $[n]\setminus \cup_{j=1}^{i-1}A_j$ zu $A_i$ gehört}
        \[ 
            \mathbb{P}(\Pi_n = \{ A_1 ... A_k\}) = p(|A_1|, ... |A_k|)
        \]  
        für eine Funktion p von Kompositionen $\textbf{n}:= (n_1, ..., n_k)$ von $n$ gilt. Diese Funktion wird auch als PEPF (Partially Exchangeable Probability Function) bezeichnet.
        \item Falls die Funktion $p$ zusätzlich noch symmetrisch ist, so ist die zufällige Partition $\Pi_n$ Austauschbar und man bezeichnet $p$ als \textit{EPPF} (Exchangeable Partition Probability Function).
    \end{enumerate}
\end{Definition}
\begin{Bemerkung}
    Das bei einer symmetrischen Funktion $p$, die zufällige Partition tatsächlich Austauschbar ist, folgt aus der Definition und der Diskussion in \cite[Seite 85]{aldous2006ecole}: eine Partition heißt austauschbar, falls 
    \[
         \mathcal{L}(\Pi_n) = \mathcal{L}(\pi(\Pi_{n})),
    \]
    für jede Permutation $\pi$ von $\{1...n\}$ gilt. Hier wirkt $\pi$ auf Partitionen $(A_1,...A_k)$ von $n$ über $\pi(A_1,...,A_k) =(\pi(A_1),...,\pi(A_k)) $ und $\pi(A_i) = \{\pi(i) | i \in A_i\}$.
\end{Bemerkung}
\begin{Definition}{Konsistenz}\\
    Wir nennen eine Folge von austauschbaren zufälligen Partitionen $(\Pi_n)$ \textit{konsistent}, falls die Restriktion $\Pi_{m,n}$ von $\Pi_m$ zu $[n]$ für alle $m > n$ die gleiche Verteilung hat wie $\Pi_m$. Alle zufälligen Partitionen auf $\mathbb{N}$ sind konsistent.
\end{Definition}
\begin{lemma}
    \label{lemma partielle austauschbarkeit}
    Sei $\Pi_n$ eine partiell austauschbare zufällige Partition von $[n]$ mit \textit{PEPF} $p(\textbf{n})$ und $\Pi_m$ die Restriktion von $\Pi_n$ auf $[m], m \in \{1...n\}$, dann gilt mit $\textbf{m} := (m_1,...,m_k)\in \bigcup_{i = 1 }^{n} \mathbb{N}^i$, $\sum_{i=1}^{k}m_i = m$ und $\textbf{m}^{i+} = \textbf{m} + e_i, i \in \{1...k+1\}$ für einen Einheitsvektor $e_i$:
\begin{center} 
    $\Pi_m$ ist partiell Austauschbar mit PEPF $p(\textbf{m}):=  \sum_{j=1}^{k+1} p(\textbf{m}^{j+})$.
\end{center}
\end{lemma}
\begin{proof}
    Leichte Induktion, siehe \cite[Proposition 10]{pitman1995exchangeable}
\end{proof}
\begin{Bemerkung}
    Es wird sich später erweisen, dass wir für bestimmte Konstruktionen von austauschbaren Partitionen von $\mathbb{N}$ die \textit{EPPF} explizit nachprüfen sollen. Dafür ist folgendes Lemma nützlich.
\end{Bemerkung}
\begin{lemma}
    \label{Austauschbar Lemma}
    Sei $N_m$, $1 \leq m \leq n$ ein zufälliges Element in der Größe der Klassen einer Partition $\Pi_m$ in Reihenfolge ihres Erscheinens. Falls $(N_1,...,N_n)$ eine Markov Kette ist mit 
    \begin{equation}
    \label{markov chain}
    \mathbb{P}(N_{m+1} = \textbf{n}^{i+} | N_{m} = \textbf{n}) = \frac{p(\textbf{n}^{i+})}{p(\textbf{n})} 
    \end{equation}
   für ein $p: \bigcup_{i = 1}^{n}\mathbb{N}^i \rightarrow [0,1]$,$p(1) = 1$, dann ist $\Pi_n$ partiell Austauschbar mit \textit{PEPF} $p$.
\end{lemma}
\begin{proof}
    Wir bemerken zuerst aus Definition \ref{Austauschbarkeit}, falls wir $N^*_n:= \{\textbf{n}:= (n_1,...,n_k)|k \in \mathbb{N},\sum_{i=1}^{k}n_i = n\}$ definieren, dass $p$ eine \textit{PEPF} ist, falls $S_n := \sum_{\textbf{n} \in N^*_n}\#(\textbf{n})p(\textbf{n}) = 1$ und $p(\textbf{n}) \in [0,1]$. Hier ist 
    \begin{equation}
        \label{combinatorial equation}
        \#(\textbf{n}):= \frac{n!}{n_k(n_k + n_{k-1})...(n_k+...+n_1)\prod_{i=1}^{k}(n_i-1)!},
    \end{equation}
    die Anzahl der Partitionen von $[n]$ mit Klassengrößen in Reihenfolge ihres Erscheinens durch $\textbf{n}$ gegeben (Siehe Bemerkung \ref{combinatorial bemerkung}). Angenommen $S_n \neq 1$, dann gilt aber $S_n> 0$ und wir definieren $q(\textbf{n}) := \frac{p(\textbf{n})}{S_n}$. Dann ist $q(\textbf{n})$ aus Lemma \ref{lemma partielle austauschbarkeit} zusammen mit \ref{markov chain} eine \textit{PEPF}, also $q(1) = \frac{p(1)}{S_n} = 1 \Rightarrow S_n =1$, somit ist $p(\textbf{n})$ eine \textit{EPPF} und die Behauptung ist gezeigt. 
\end{proof}
\begin{Bemerkung}
    \label{combinatorial bemerkung}
    Das $\#(\textbf{n})$ tatsächlich die Form \ref{combinatorial equation} hat, wird im Anhang, Lemma \ref{combinatorial lemma}, bewiesen und ist für die restliche Ausarbeitung nicht weiter relevant.
\end{Bemerkung}
\end{section}
\begin{section}{Chinese Restaurant Process}
% LTeX: language=de-DE
    Der \glqq Chinese Restaurant Process\grqq, dient der Konstruktion von zufälligen austauschbaren Partitionen, mithilfe welcher die Analyse von Bäumen erleichtert wird. In Kapitel \textcolor{red}{add something}, wird eine Bijektion zwischen den Bäumen und Partitionen aufgestellt und dies erlaubt uns, die in diesem Kapitel entwickelte Theorie anzuwenden. Wir beginnen mit der Konstruktion vom \glqq Chinese Restaurant Process\grqq.\\
\newline
\textbf{\fontsize{14}{18}\selectfont Konstruktion von CRP}
\\
Bevor die Konstruktion vom Prozess beschrieben wird, wollen wir vorab das Ziel der Konstruktion erläutern. Die Konstruktion erlaubt uns eine konsistente Folge von austauschbaren zufälligen Partitionen von $[n]$ zu erzeugen, also eine zufällige Partition $\Pi_\infty := (\Pi_n)_{n \in \mathbb{N}}$ von $\mathbb{N}$ zu generieren. Dies erlaubt uns u.a. Kingmans Sätze wie \textit{Kingman's Correspondance} und \textit{Kingman's representation} zu nutzen und somit Aussagen über die Grenzwerte der Größe der Blöcke von Partitionen zu formulieren. Wir beginnen mit einer einfachen Form des Prozesses: \\
\\
Wir betrachten eine Folge von zufälligen Permutationen $ \{\sigma_n, n \in \mathbb{N}\}$, sodass 
\begin{enumerate}
    \item $ \{\sigma_n\}$ Uniform verteilt auf $[n]$ ist.
    \item $\forall n \in \mathbb{N}$, falls $\sigma_n$ als Produkt von Zyklen dargestellt wird, $\sigma_{n-1}$ von $\sigma_n$ durch Deletion von Element $n$ zurückgewonnen wird.   
\end{enumerate}
Die Verteilung, die durch diese zwei Regeln entsteht, kann durch einen zufälligen Sitzplan in einem chinesischen Restaurant beschrieben werden. Man stelle sich das Restaurant zunächst vor als eine Sammlung von abzählbar unendlichen vielen Tischen, wo jeder Tisch die Sitzplatzkapazität von abzählbar unendlich vielen Kunden hat. Die Kunden, nummeriert mit den natürlichen Zahlen nach Reihenfolge Ihres Erscheinens, werden nach folgenden Regeln platziert.
\begin{enumerate}
    \item Der erste Kunde sitzt an Tisch 1
    \item Falls zum Zeitpunkt des Eintritts des $i$-ten Kundes $k$ Tische besetzt sind, so setzt sich dieser mit gleicher Wahrscheinlichkeit $(\frac{1}{n+1})$ links von Kunden $i: i \in \{1...n\}$ oder an einem neuen Tisch $k+1$.
\end{enumerate}
Wir definieren nun $\sigma_n: [n] \rightarrow [n]$ mithilfe dieser Regeln, sodass falls Kunde $i$ links von Kunde $j$ sitzt, $\sigma_n(i) = j$ und falls Kunde $i$ alleine sitzt $\sigma_n(i) = i$. Dass diese Folge von
Permutationen die ersten 2 Regeln befolgt, ist durch eine triviale Induktion ersichtlich. 
\\
Die Folge, die durch die Zyklen von den Permutationen generiert werden, ist austauschbar und konsistent, generiert somit eine austauschbare Folge $\Pi_\infty$ von $\mathbb{N}$.
    %\item Falls zum Zeitpunkt des Eintritts des $i$-ten Kundes $k$ Tische besetzt sind, so setzt sich dieser an Tisch $j \in {1...k}$ mit Wahrscheinlichkeit $\frac{n_i-\alpha}{n+\theta}$ und an Tisch $k+1$ mit Wahrscheinlichkeit $\frac{\theta + k\alpha}{n + \theta}$
\end{section}
\begin{section}{Verallgemeinerte CRP}
Die einfache Konstruktion vom \textit{Chinese Restaurant Process} lässt sich recht einfach verallgemeinern, sodass eine größere Klasse von austauschbaren Partitionen betrachtet werden können. Anstatt, wie im einfachen Fall, nur uniform ein Sitzplatz auszuwählen (Regel 2), wird in der allgemeinen Form des Prozesses eine beliebige Folge $(P_1,P_2...)$ von Zufallsvariablen, mit $\sum_{i\geq1}P_i \leq 1$ und $P_i \geq 0$ gewählt, und die zweite Regel mit folgender ersetzt:
\begin{itemize} 
    \item Falls zum Zeitpunkt des Eintritts des i-ten Kundes $k$ Tische besetzt sind, so setzt sich dieser mit Wahrscheinlichkeit $P_j, j \in \{1...k\}$ an Tisch $j$ und mit Wahrscheinlichkeit $1 - \sum_{j=1}^{k}P_j$ an einen neuen Tisch.
\end{itemize}
Dieser Prozess erlaubt es nicht nur jede austauschbare Folge zu konstruieren, sondern tatsächlich jede partiell austauschbare Folge zu konstruieren. Beruhend auf dieser Konstruktion, hat Pitman \cite{pitman1995exchangeable} folgende Variation von \textit{Kingman's Representation} bewiesen
\begin{theorem}
    \label{pitman}
    Sei $\Pi_\infty$ eine zufällige Partition von $\mathbb{N}$, dann sind folgende Äquivalent
    \begin{enumerate}
        \item $\Pi_\infty$ ist partiell Austauschbar
        \item Falls wir die Tische mit Partitionen identifizieren, dann wird $\Pi_\infty$ gegeben $(P_1,P_2...)$ durch den verallgemeinerten \textit{Chinese Restaurant Process} generiert. Also $\forall{n} \in \mathbb{N}$ bedingt gegeben $(P_1,P_2...)$ und $\Pi_n = \{A_1,A_2,...A_k\}$, mit $A_i, i \in \{1...k\}$ in Reihenfolge ihres Erscheinens, $\Pi_{n+1}$ eine Erweiterung von $\Pi_n$ ist, wobei $n+1$ zur Klasse $A_i$ mit Wahrscheinlichkeit $P_i$ angefügt wird und zur einer neuen Klasse $A_{k+1}$ mit Wahrscheinlichkeit $1 - \sum_{i=1}^{k}P_i$ angefügt wird. 
    \end{enumerate}
Weiter sind die $P_i$ in 2. fast sicher eindeutig und gleich der fast sicheren asymptotischen Frequenz der Klassen $A_i$, also
\[ 
\frac{|A_i \cap [n]|}{n} \to P_i \hspace{10pt} \text{P -f.s.} \hspace{5pt} \text{für } n \to \infty.
\]
\end{theorem}
\begin{proof}
    Dieser Satz wird in \cite{pitman1995exchangeable} bewiesen.
\end{proof}
\begin{Bemerkung}
    Der Vorteil von Pitmans Charakterisierung von partiell Austauschbaren Folgen gegenüber \textit{Kingman's Representation} ist die Bestimmung der Verteilung $(P_1,P_2...)$ im Vergleich zu der sortierten Folge $(P^\downarrow_1,P^\downarrow_2,...)$, welche durch die Konstruktion vom \textit{Chinese Restaurant Process} erleichtert wird. Mithilfe vom Korollar \ref{korollar pitman} lässt sich die $PEPF$, bzw. $EPPF$ im austauschbaren Fall, recht einfach berechnen.
\end{Bemerkung}
\begin{Korollar}
    \label{korollar pitman}
    Die PEPF $p$ von Satz \ref{pitman} hat Form
    \begin{equation}
    \label{PEPF expected value}    
    p(n_1,n_2,...,n_k) = \mathbb{E}\left[\prod_{i=1}^{k}P_i^{n_i-1}\prod_{i=1}^{k-1}(1-\sum_{j=1}^{i}P_j)\right].
    \end{equation}
\end{Korollar}
\begin{proof}
    Nach Definition \ref{Austauschbarkeit}, gilt mit $\mathbb{A}_n:= \{A_1,...A_k\}$ und $\mathbb{A}_{n-1}:= \mathbb{A}_n \cap [n-1]$ und Satz \ref{pitman}
    \begin{align*}
    p(|A_1|,...,|A_k|) &= \mathbb{P}(\Pi_n = \mathbb{A}) \\
    &= \mathbb{E}\left[\mathbb{P}(\Pi_n = \mathbb{A}_n|(P_1,P_2...))\right]  \\
    &= \mathbb{E}\left[\mathbb{P}(\Pi_n = \mathbb{A}_n|\Pi_{n-1} = \mathbb{A}_{n-1},(P_1,P_2...))\mathbb{P}(\Pi_{n-1} = \mathbb{A}_{n-1}|(P_1,P_2...))\right] \\
    &=... =  \mathbb{E}\left[\prod_{i=1}^{k}P_i^{n_i-1}\prod_{i=1}^{k-1}(1-\sum_{j=1}^{i}P_j)\right].
\end{align*}
\end{proof}
Mithilfe einer anderen Konstruktion von austauschbaren $\mathbb{N}$ \cite[Konstruktion 16]{pitman1995exchangeable}, stellt sich heraus, dass die $P_i$ in Form von einer \textit{Residual allocation Model (RAM)} sich darstellen lassen. Das bedeutet, dass unabhängige Zufallsvariablen $W_i, i \geq 1, W_i \in [0,1]$ existieren, sodass die $P_i$ aus Satz \ref{pitman} als Produkt von den $W_i$
\begin{equation}
    \label{RAM probability}
    P_i = W_i\prod_{j=1}^{i-1} (1-W_j)
\end{equation}
darstellen lassen. Korollar \ref{korollar pitman} ändert sich somit zu
\begin{Korollar}
    \label{second korollar Pitman}
    Die PEPF $p$ von Satz \ref{pitman} hat Form
    \[
    p(n_1,n_2,...,n_k) = \mathbb{E}\left[\prod_{i=1}^{k}W_i^{n_i-1}(1-W_i)^{\sum_{j=i+1}^{k}n_j}\right].
    \]
\end{Korollar}
\begin{proof}
   Wir nutzen \ref{RAM probability} und Induktion um zuerst
   \[
    \prod_{i=1}^{n}(1-W_i) = 1 - \sum_{i=1}^{n}P_i
   \] 
   zu zeigen.\\
   $n = 1$: klar. \\
   Sei die Aussage für $n$ Wahr, dann gilt 
   \begin{align*}
    1 - \sum_{i=1}^{n+1}P_i &\overset{I.A.}{=} \prod_{i=1}^{n}(1-W_i) - P_{n+1}\\
    &=  \prod_{i=1}^{n}(1-W_i) - W_{n+1}\prod_{i=1}^{n}(1-W_i) \\
    &= \prod_{i=1}^{n+1}(1-W_i).
   \end{align*}
   Eingesetzt in \ref{PEPF expected value}, liefert 
   \begin{align*}
    p(n_1,n_2,...,n_k) &= \mathbb{E}\left[\prod_{i=1}^{k}W_i^{n_i-1}\prod_{j=1}^{i-1}(1-W_j)^{n_i-1}\prod_{i=1}^{k-1}\prod_{j=1}^{i}(1-W_j)\right]\\
    &= \mathbb{E}\left[\prod_{i=1}^{k}W_i^{n_i-1}\prod_{j=1}^{i-1}(1-W_j)^{n_i-1}\prod_{i=1}^{k}\prod_{j=1}^{i-1}(1-W_j)\right]\\
    &= \mathbb{E}\left[\prod_{i=1}^{k}W_i^{n_i-1}\prod_{j=1}^{i-1}(1-W_j)^{n_i}\right]\\
    &=\mathbb{E}\left[\prod_{i=1}^{k}W_i^{n_i-1}(1-W_i)^{\sum_{j=i+1}^{k}n_j}\right].
   \end{align*}  

\end{proof}
Wir nutzen die bis jetzt entwickelte Theorie der austauschbaren Partitionen um eine Klasse von \textit{Chinese Restaurant Process} zu betrachten, welche für spätere Baumanalyse von Relevanz wird. Dafür ersetzen wir die zweite Regel in der Konstruktion vom verallgemeinerten \textit{Chinese Restaurant Process} mit genauen Wahrscheinlichkeiten. Für Parameter $\alpha$ und $\theta$ halten wir fest:
\begin{itemize} 
    \item Falls zum Zeitpunkt des Eintritts des i-ten Kundes $k$ Tische mit jeweils $(n_1,...,n_k)$ Kunden besetzt sind, so setzt sich dieser mit Wahrscheinlichkeit $\frac{n_j-\alpha}{n + \theta}$ an Tisch $j \in \{1...k\}$ und mit Wahrscheinlichkeit $\frac{\theta + k\alpha}{n + \theta}$ an einen neuen Tisch $k+1$.
\end{itemize}
Damit die Wahrscheinlichkeitsregeln bei dieser Konstruktion erhalten bleiben, müssen $\alpha$ und $\theta$ entsprechend gewählt werden, dafür gibt es 2 Fälle \textcolor{red}{add proof of this}: 
\begin{enumerate}
    \item $ 0 \leq \alpha < 1$ und $\theta > - \alpha $ 
    \item $\alpha < 0$ und $\theta = -m \alpha$ für ein $m \in \mathbb{N}$
\end{enumerate}

Das Haupttheorem ist nun folgendes.
\begin{theorem}
    \label{Main theorem CRP}
    Für jede der Parameter $(\alpha,\theta)$, welche die oben genannten Voraussetzungen erfüllen, generiert das CRP eine austauschbare  Partition von $\mathbb{N}$. Die dazu gehörende EPPF (siehe Def. 3) ist 
    \begin{equation}
   p_{\alpha,\theta}(n_1,...,n_k) = \frac{\displaystyle \prod_{i=0}^{k-2}(\theta + \alpha + i\alpha)\prod_{j=1}^{k}\prod_{r=1}^{n_j-1}(r-\alpha)}{\displaystyle\prod_{i=1}^{n-1}(\theta  + i)}.
    \end{equation}
    Die zugehörigen asymptotischen Frequenzen, in \textit{Reihenfolge ihres Erscheinens}, können folgendermaßen repräsentiert werden.
    \[
    (P_1,P_2,P_3...) = (W_1,\bar{W_1}W_2,\bar{W_1}\bar{W_2}W_3...).
    \]
    Wobei $W_i$ unabhängig voneinander, nach $\beta(1 - \alpha, \theta + i\alpha)$ verteilt sind und $\bar{W_i}:= 1 - W_i$.
\end{theorem} 
Bevor wir zum Beweis kommen, definieren wir die $\beta$-Verteilung und nennen Eigenschaften.
\begin{Definition}
    Eine Zufallsvariable $X$ ist $\beta(a,b),$ mit $ a,b > 0 $ verteilt, falls $X$ die Wahrscheinlichkeitsdichte 
    \[
        f(x) =
        \begin{cases}
        \frac{1}{B(a,b)} x^{a-1}(1-x)^{b}, \hspace{5pt} &x \in [0,1]  \\
        0 &\textnormal{sonst},
        \end{cases} 
   \]
mit 
\[
    B(a,b):= \frac{\Gamma(a+b)}{\Gamma(a)\Gamma(b)}
\]
besitzt. Hier bezeichnet $\Gamma$ die Gamma Funktion, $\Gamma(z)= \int_{0}^{\infty} t^{z-1}e^{-t}dt$.
\end{Definition}
Aus Korollar \ref{second korollar Pitman} werden oft nur die Momente von $W_i$ zur Verfügung stehen, woraus die Verteilung abgeleitet werden soll. Für Zufallsvariablen mit beschränkten Träger ist die Verteilung durch die Momente festgelegt, \cite[Theorem 30.1]{billingsley1995probability}, insbesondere auch für die $\beta$-Verteilung.
\begin{lemma}
    Für $X \sim \beta(a,b)$, gilt
    \[
    \mathbb{E}[X^k] = \prod_{r=0}^{k-1}\frac{a+r}{a+b  + r}.
    \]
\end{lemma}
\begin{proof}
    Die Momente werden über die Berechnung der momenterzeugenden Funktion von $X$ in \cite[Kapitel 25]{johnson1995continuous} bewiesen. 
\end{proof}

Insbesondere folgt für die Parameter aus Satz \ref{Main theorem CRP}: 
\begin{Korollar}
    \label{Korollar Momente beta}
    Für $X \sim \beta(1-\alpha,\theta + i\alpha)$, gilt:
    \[
    \mathbb{E}[X^k] = \prod_{r=1}^{k}\frac{r-\alpha}{r+\theta +(i-1)\alpha }.
    \]
\end{Korollar}
Mit diesem Hilfsmittel können wir nun den Satz beweisen.
\begin{proof}[Beweis von Satz \ref{Main theorem CRP}]
   Wir beginnen mit der Austauschbarkeit. Dafür prüfen wir zuerst, dass $p_{\alpha,\theta}$ tatsächlich unserer EPPF entspricht. Wir nutzen Lemma \ref{Austauschbar Lemma} und beachten, dass $p(1) = 1$. Weiter sind die Übergangswahrscheinlichkeiten von $p_{\alpha,\theta}$ genau die Übergangswahrscheinlichkeiten von $N_m$ aus dem Lemma. Nach Definition \ref{Austauschbarkeit} ist die \textit{PEPF} $p$ symmetrisch, also ist die Konstruktion austauschbar. Für die asymptotischen Frequenzen nutzen wir Lemma \ref{second korollar Pitman} und verfahren Induktiv um die $W_i$ zu bestimmen. Wir beachten, dass die Momente von $W_i$ eindeutig die $W_i$ charakterisieren, da Sie nach Voraussetzung beschränkten Träger in $[0,1]$ haben. Für $W_1$ sind die Momente nun:
    \begin{align*}
    p(i)= \mathbb{E}[W_1^{i-1}] &= \frac{\prod_{j=1}^{i-1}(j-\alpha)}{\prod_{j=1}^{i-1}(\theta+j)}\\
    &= \prod_{j=1}^{i-1}\frac{(j-\alpha)}{(\theta+j)}
    \end{align*}
und somit $W_1 \sim \beta(1-\alpha,\theta + \alpha)$ nach Korollar \ref{Korollar Momente beta}. Wir machen eine Induktion über $W_i$ und nehmen an, dass $W_n \sim \beta(1-\alpha,\theta+n\alpha)$. Für $W_{n+1}$ rechnen wir
    \begin{align*}
    \frac{p(\overbrace{1,..,1}^{n},i+1)}{p(\underbrace{1,..,1}_{n-1},i+2)}= \frac{\mathbb{E}[(1-W_n)^{i+1}]}{\mathbb{E}[(W_n)^{i+1}]}\mathbb{E}[W_{n+1}^i] \\
    \Rightarrow \frac{\theta + n\alpha}{i+1-\alpha}=  \prod_{r=0}^{i}\frac{\theta + n\alpha + r}{{1-\alpha + r}}\mathbb{E}[W_{n+1}^i]\\
    \Rightarrow \frac{\theta + n\alpha}{i+1-\alpha}\prod_{r=1}^{i+1}\frac{{r-\alpha }}{ r -1 + \theta + n\alpha }=  \mathbb{E}[W_{n+1}^i]\\
    \Rightarrow \prod_{r=2}^{i+1}\frac{1}{r -1 + \theta + n\alpha}\prod_{r=1}^{i}r-\alpha =  \mathbb{E}[W_{n+1}^i]\\
    \Rightarrow \prod_{r=1}^{i}\frac{r-\alpha}{r+\theta+n\alpha} = \mathbb{E}[W_{n+1}^i].
    \end{align*}
Die Aussage folgt.
\end{proof}
\begin{Definition}
Die Verteilung der $(P_1,P_2...)$ aus Satz \ref{Main theorem CRP} wird nach Griffith-Engen-McCloskey für Parameter $\alpha,\theta$ die $GEM(\alpha,\theta)$-Verteilung genannt.
\end{Definition}
\end{section}
\chapter{Linear preferential trees} 
% LTeX: language=de-DE
%Das Kapitel beschreibt den Aufbau von \text{linear preferential trees} und wie diese zu Split Bäume 
In diesem Kapitel werden \textit{general preferential trees} eingeführt und die Einbettung von diesen in Devroyes \textit{Split} Bäume besprochen. Wir beginnen mit der Konstruktion von den jeweiligen Bäumen.
\\
\\
\textbf{\fontsize{14}{18}\selectfont Konstruktion von general preferential trees}\\
Der Aufbau von \textit{general preferential attachment trees} folgt im Wesentlichen dem Aufbau von \textit{random recursive Trees}. Wobei in \textit{Random recursive Trees} Knoten aufeinanderfolgend, mit uniformer Wahrscheinlichkeit als Kind von bereits existierenden Knoten angefügt werden, ist das Anknüpfen von Knoten von \textit{general preferential attachment trees} noch abhängig vom Ausgangsgrad von den existierenden Knoten. Präziser formuliert, wird der Vaterknoten $v$ von einem neuen Knoten, proportional zu $w_d(v)$ ausgesucht, wobei $d(v)$ der Ausgangsgrad vom Knoten $v$ ist, und $w_1, w_2...$ eine Folge von Gewichten ist. Der Fall $w_i = 1, \forall i \in \mathbb{N}$, führt somit genau auf den Fall von \textit{random recursive trees} zurück. In dieser Ausarbeitung werden die Gewichte von der Form 
\[
w_k = \chi k + \rho \hspace{5pt}, \chi \in \mathbb{R}\hspace{5pt} \rho > 0
\]
betrachtet, woher der Name \textit{linear preferential attachment trees} auch stammt. Da der Vaterknoten lediglich proportional zu der Folge von Gewichten ausgesucht wird, ändert die Skalierung der Folge nicht die Generierung der Baumstruktur, also beschränken wir die Betrachtung von den Gewichten auf 3 Spezialfälle:
\begin{enumerate}
    \item $\chi = 1$ ist der Fall vom \textit{general plane oriented recursive trees} \cite{panholzer2007level} 
    \item $\chi = 0$ ist der Fall vom \textit{random recursive trees}
    \item $\chi = -1$ ist der Fall von \textit{m-ary increasing trees}
\end{enumerate}
\begin{theorem}
\cite[Lemma 3.1]{janson2019random}
   Mit den linearen Gewichten $w_k = \chi k + p $, hat ein Baum $T$ mit $m$ Knoten ein Gesamtgewicht $w(T) = (m-1)\chi + m\rho = m(\chi + \rho) - \chi$.
\end{theorem}
\begin{proof}
    Seien $d_1...d_m$ die Ausgangsgrade der jeweiligen Knoten im Baum. Dann ist die Summe der Ausgangsgrade, $\sum_{i=1}^{m}d_i = m-1$ und somit:
    \[
     w(T) = \sum_{i=1}^{m}w_i = \sum_{i=1}^{m} (\chi d_i + \rho) = m \rho + \chi \sum_{i = 1}^{m} d_i = m(\chi + \rho) - \chi  
    \]
\end{proof}
\begin{theorem} \cite[Theorem 3.2]{janson2019random}
    Sei $(T_n)_1^\infty =(T_n^{\chi,\rho})_1^\infty = $ eine Folge von \linpreft, mit den Kindern von der Wurzel in der Reihenfolge ihres Erscheinens beschriftet. Sei $N_j(n) := |T^j_n|$, die Größe des j-ten \PsubT von $T_n$. Dann gilt 
    \[ 
        \forall j \geq 1  \hspace{6pt} \frac{N_j(n)}{n}  \rightarrow P_j \hspace{6pt} \text{für} \hspace{6pt} n \rightarrow \infty \hspace{6pt} \text{P-f.s.}.
  \]
Wobei die Zufallsvariablen  $P_j$ Verteilung $GEM(\frac{\chi}{\chi+\rho},{\frac{\rho}{\chi+\rho}})$ haben. \textcolor{red}{add case 0} 
\end{theorem}
\begin{proof}
    Da wir im Allgemeinen $\chi + \rho > 0 $ annehmen dürfen, und da die Multiplikation von allen $w_k$ mit einer Konstante nicht den Baum ändert, können wir o.B.d.A   $\chi + \rho = 1 $ annehmen. Das Argument folgt nun durch Übersetzung von \PsubT zu Zyklen in Permutationen und eine Bijektion zwischen \linpreft und Permutationen. Wir beginnen mit der Bijektion. \\
    Dafür beschriften wir die Knoten sukzessiv, beginnend mit der Wurzel als 0, nach der Reihenfolge vom Hinzufügen zum Baum (der $i$-ter Knoten der hinzugefügt wird, wird mit $i-1$ beschriftet). Die Wurzel, die mit 0 beschriftet wird, wird identifiziert mit der leeren Permutation und der Zweite Knoten, welcher mit 1 beschriftet wird, wird mit dem Zyklus $(1)$ identifiziert. Wir identifizieren nun unsere Bäume mit Permutationen induktiv, nach folgenden zwei Regeln:
    \begin{enumerate}
        \item Falls der Knoten mit Beschriftung $i$ an die Wurzel angehängt wird, so fügen wir zur bestehenden Zyklus-Dekomposition, die Fixpunkt-Permutation $(i)$ hinzu.
        \item Falls der Knoten mit Beschriftung $i$ an ein anderen Knoten mit Beschriftung $j$, $j \neq 0$ angehängt wird, so fügen wir in den Zyklus $C$, welcher $j$ enthält, $i$ direkt nach $j$ hinzu. Bedeutet also, falls im ursprünglichen Zyklus $j \rightarrow k$ abgebildet wurde, im nächsten Schritt $j \rightarrow i \rightarrow k$ abgebildet wird.
    \end{enumerate}
Es ist einfach zu zeigen, dass diese Abbildung von Bäumen zu Permutation, welche durch diese Konstruktion entsteht, bijektiv ist und somit können wir die Theorie die für 
\end{proof}

%\section{Einleitung}
\bibliographystyle{plain}
\bibliography{Sources/Stochastik_Literatur}
\end{document}